\chapter*{Introduction to research}
While taking on problems and resolving theories are right and interesting, the fact remains that this book's prerequisite to said conjectures indict that we need to write them into papers. That is right, papers for research publication. Which, by all means is a great format, and also there are problems of which I can only resolve or present, using it as a medium. It is just that it is fairly disorganized, and I don't actually want to do it in the majority of the time, and it might just as be as reasonable as anything can take of. Most of the time though, this branches out from my observation that the waitlist toward which a paper is submitted and accepted, is fairly long. And in due time, no can wait for my graduation. 

Those researches range from theoretical papers and inquiries on topics, to experimental results and tutorials, to expository studies on various topics and their solutions, and more on that matters toward novel ideas and theoretical treatment, almost new, and almost entirely different. Historical accounts also play a role in such list, and we will refer to them in due time. For those in the shortlist \footnote{A shortlist refers to a series of short span, short-term research papers or anything of interest. Specifically, can be used as reference papers or constructing further \textit{profile paper} (used in CV and academic records).}, they would be recorded differently, and will be held within a certain standard. 

The research part of this book will contain several components related to its organization, and definitely its content in serving. First, here, is the introductory session on what the researches are about, logs about them, records, and organization notes - complete, incoming, stalled, delayed, scrapped papers, et cetera. Then the main contents, and finally, supplementary materials and thereof. Not too much to be of concern, but rather cumbersome to manage. Each research paper then belongs to a certain chapter-wise, so expect the table of content to be ridiculously long in part. 

\section{Research topics and organization}

Because of the lacking of a corrected list of research, and the lack of also resources to run most of them, leaving the research topic, of which in fair spirit can be hard to come by again, out alone, is perhaps not so ideal from a given perspective. Thereby, one of the main goal of this section is then to refute that, and coming together with a list of reearch proposals and topics that receive fairly good assessment, coupled together, and filed into a singular space. 

\section{Research entries}

Starting from here, we have the entries of paper ideas and titles, established topics and thereof. Some of them are already included, some of them are not, but will soon be included into the list. There are two main types of research entries - the first type is the \textbf{theoretical entries} of theory papers and thereof, analysis, examination and experimental results, confirmation and investigation of different phenomena. in general, an \textit{academically close} type of entry. The second one is the \textbf{construction entries}, of implementation, structural formation, building an actual system or making a workable framework, developing a principle for example, and so on. While this is not fully rigorous in category, for organizing template, it is good enough. 

\begin{enumerate}
    \item \textbf{Double descent: a review} (review, reference): A reference and review paper talking about double descent. Including effort to formalize the definition of it. 
    \item \textbf{Interpreting double descent in polynomial via density estimation} (original, theory): A fairly light and theoretical paper on the toy model description of polynomial model and different interpretation. 
    \item \textbf{Double descent on support vector machine and polynomial models on binary classification - an analysis} (experimental, theory): A both theory and experimental paper. 
    \item \textbf{PINNs and double descent} (experimental): An experimental on PINNs and how to identify double descent on it. 
    \item \textbf{A review of model structures in machine learning and theoretical machine learning} (theoretical, review, reference): A reference paper on the concept of models in machine learning, their structures, flavour, consideration, and thereof. 
    \item \textbf{Different treatment of machine learning - an analysis} (review, reference, original): A paper on the different treatment of machine learning, from different lens and scale, and how do they fit together (categorical machine learning for example). 
    \item \textbf{Reuniting the neural network framework using partial category theory} (original, theory): Using the proposed categorical machine learning or neural network to consider reuniting different architecture together under the same framework. 
    \item \textbf{On the capacity and capability of neural networks} (review, original): A paper analysing the capacity and capability of neural networks on various tasks and setting thereof, and so forth, considering different structures. 
    \item \textbf{Machine learning from a mathematical modelling view} (review, original): Proposing a modelling theoretic, internal state interpretation to machine learning. 
    \item \textbf{The analysis of statistical physics on machine learning} (review, original): Setting up a review work on how statistical physics is related to machine learning, where is it failing and what constitute the failure. 
    \item \textbf{Classical model into neural network architecture - an analysis} (original, theoretical): An attempt to attack and formalize classical models into neural network theoretic (unit-wise approach of neuronal unit) to classical models, hence gauging their capability. 
    \item \textbf{Unsolved problems in theoretical machine learning} (original, review): A work in review of different interpretations and problems residing the current framework of machine learning. 
    \item \textbf{Differential equations and machine learning - interpreting machine learning system using differential equations} (original, theory): An attempt to express ML systems in terms of differential equations, just as the differential model. 
    \item \textbf{The theory of Raman Spectroscopy} (review, reference): On the theory of Raman spectroscopy on itself, different interpretation, history, approaches, and problems. 
    \item \textbf{Permuting layer optimization of neural network}: A proposal on which in addition to optimize by weights, we optimize the neural network model also by its own layers. 
    \item \textbf{Non-layer analysis of neural network} (experimental, theory): An analysis of neural network architecture and formalism, using a non-layer interpretation and experimental method, with the tradeoff of reduced computational power. 
    \item \textbf{A polynomial toy model analysis of double descent} (theoretical, original): An analysis on polynomial model approximation and double descent on said model. 
    \item \textbf{Raman spectroscopy and machine learning analysis of knowledge-aware model} (theory, experimental, original): Work on analysis and experimental construction of Raman spectroscopy models, with knowledge of the system instead of relying on only data approximation. 
    \item \textbf{Model complexity and representation space - an old idea analysis} (original, theory): An inquiry on the theory of model complexity and representational complexity typically seen in machine learning theory. 
\end{enumerate}

There are topics marked as \textit{original}, meaning a major portion of that paper is original and not recreation or provision from other papers; \textit{experimental} as in with experimental results and insight analysis, \textit{theory} as in theoretical treatment and formalism, and finally \textit{review} and \textit{reference} as in expository-wise and exploratory papers. 

\subsection{Constructions}
Construction is well, construct instruction so to speak. Essentially, while it is good to just experiment, it is also detrimental to develop based on such. Truth is, we need a nice ROE on research to practical application and non-toy analysis. 


